%%%%%%%%%%%%%%%%%%%%%%%%%%%%%%%%%%%%%%%%%%%%%%%%%%%%%
%			HLAVIČKA								%
%%%%%%%%%%%%%%%%%%%%%%%%%%%%%%%%%%%%%%%%%%%%%%%%%%%%%
\documentclass[openany,12pt]{memoir}

\usepackage[utf8]{inputenc} 
\usepackage[czech]{babel}
\usepackage[T1]{fontenc}
\usepackage[top=1.5cm, bottom=2cm, left=2cm, right=2cm]{geometry}  % --> NASTAVENÍ OKRAJŮ
\usepackage{fancyhdr}
\usepackage{graphicx}
%\usepackage{lmodern}
\usepackage{tgschola}  % --- Silnější font
\usepackage{xwatermark}
\usepackage{xcolor}
\usepackage{changepage}
\usepackage{newunicodechar}



%%%%%% Package na zpěvník
\usepackage[full]{leadsheets}%http://mirrors.nic.cz/tex-archive/macros/latex/contrib/leadsheets/leadsheets_en.pdf   --> dokumentace	
\definesongtitletemplate{empty}{} 
\setchords{
format = \bfseries,   %tučné akordy
minor = {mi},% 
input-notation = {german},%
output-notation = {german}%
}
\definesongtitletemplate{empty}{} 

\newlength{\drop}
\newwatermark[allpages,color=red!50,angle=0,scale=2, xpos=0,ypos=0]{\includegraphics[width=5cm]{obr/pozadi2.jpg}} %--> dvojka na pozadí


%%%% Vlastní příkazy
\newcounter{Slokočet}   %Automatické číslování slok
\newcommand{\mezera}{\vspace*{0.5cm}}   %Horizontální odsazení slok
\newcommand{\stred}{5.2cm}   %%% Na zarovnání slok doprostřed, pozn. automatičtější zarovnávání na střed nejde
\newcommand{\refren}{\mezera \noindent \textbf{R:} } %refrén
\newcommand{\sloka}{\mezera \noindent \addtocounter{Slokočet}{1} \arabic{Slokočet}. } 	%sloka, která se automaticky čísluje
\newcommand{\ssloka}{\mezera \noindent}  % vlastní číslo sloky
\newcommand{\carka}{,\:}
\newcommand{\predehra}{\mezera \hspace*{-1.85cm} \textbf{Předehra:} }
\newcommand{\m}[1]{\color{white}{#1}}  %Pro akordy
\newcommand{\ap}{'} %Pro apostrof
\newcommand{\elipsa}{\kern\fontdimen3\font} %Příkaz pro lepší zacházení s výpustkami (=...); je to vpodstatě jen mezera mezi tečkama výpustky

%%%%%%%%%%%%%%%%%%%%%%%%%%%%%%%%%%%%%%%%%%%%%%%%%%%%%%
%			NÁVOD									 %
%%%%%%%%%%%%%%%%%%%%%%%%%%%%%%%%%%%%%%%%%%%%%%%%%%%%%%
%1. Věci v hlavičce IGNOROVAT
%2. Píseň psát do prostoru mezi \begin{song} a \end{song}
%3. další řádek se značí dvěma odsazeníma (= dvakrát stisknout enter)
%4. \refren vždy na začátku refrenu a \sloka na začátek sloky (automaticky se čísluje)
%  \ = alt gr + q ; [/] = alt gr f/g ; {/} = alt gr + b/n; ^ = alt gr + 3 + mezera
%Cokoliv napíšete do ^{  } se bude brát jako akord
%když se toto bude dotýkat nějakého slova (nebude mezi tím a slovem mezera)
%tak se akordy zjeví nad slovem, ale pište to před slovo
%Když se to nedotýká slova, tak akord lítá ve vzduchu a vytiskne se větší mezera
%První možnost je asi preferovanější
%5. Akordy stačí psát jen do první sloky, když se nezmění -- kytaristi to zvládnou



%\usepackage{subfiles}
%</++>

\begin{document}

\pagestyle{empty}
%\addtocontents{toc}{\protect\thispagestyle{empty}}
%\tableofcontents \thispagestyle{empty}\newpage
\newgeometry{top=1.5cm, bottom = 0cm, left = 2cm, right = 2cm}

\newpage

%Pro kompilaci je důležitý, 	aby ten input a \end{document} zůstali na posledních
%dvou řádcích!!!!
\begin{song}{title=\centering Veličenstvo Kat \\\normalsize Karel Kryl  \vspace*{-0.3cm}}  %% sem se napíše jméno songu a autor
\moveright -0.9cm \vbox{      %Varianta č. 1  ---> Jeden sloupec zarovnaný na střed	
\begin{minipage}[t]{0.55\textwidth}\setlength{\parindent}{0.45cm}  %Varianta č. 2 --> Dva sloupce

\sloka
	V ^{Ami{\color{white}\_\_\_}}ponurém osvětlení ^{G{\color{white}\_\_\_\_\_}}gotického ^{Ami\,\,}sálu
 
	^{C{\color{white}\_\_\_\_}}kupčíci vyděšení ^{Emi{\color{white}\_}}hledí do ^{E}misálů

	a ^{Ami{\color{white}\_\_}}houfec ^{F{\color{white}\_\_\_\_}}mordýřů si ^{G{\color{white}\_\_}}žádá ^{{\color{white}\_\_\_\_\_}C\,\,G}požehnání,
         
	/: ^{Dmi{\color{white}\_}}vždyť první z ^{Ami{\color{white}\_\_}}rytířů je ^{E7{\color{white}\_\_\_\_\_\_}}Veličenstvo ^{Ami}Kat. :/

\sloka
	Kněz-ďábel, co mši slouží, z oprátky má štolu,
	
	pod fialovou komží láhev vitriolu,
	
	pach síry z hmoždířů se valí k rudé kápi

	/: prvního z rytířů, hle: Veličenstvo Kat. :/
	
	
\refren
	^{C}Na korouhvi ^{G{\color{white}\_\_\_}}státu je ^{F{\color{white}\_\_\_\_}}emblém s ^{G}gilotinou,

	z ^{C{\color{white}\_\_\_\_\_}}ostnatýho ^{G{\color{white}\_\_\_}}drátu ^{F{\color{white}\_\_\_}}páchne to ^{{\color{white}\_}G}shnilotinou,

	v ^{Dmi\,\,}kraji hnízdí hejno ^{{\color{white}\_\_\_\_\_}Ami}krkavčí,

	^{Dmi\,\,}lidu vládne mistr ^{{\color{white}\_\_\_\_\_}E7}popravčí.

\sloka
	Král klečí před Satanem na žezlo se těší
   
	a lůza pod platanem radu moudrých věší

	a zástup kacířů se raduje a jásá,
 	
 	/: vždyť prvním z rytířů je Veličenstvo Kat. :/

\sloka
	Na rohu ulice vrah o morálce káže,
   
	před vraty věznice se procházejí stráže,
  
	z vojenských pancířů vstříc černý nápis hlásá,
   
   	/: že prvním z rytířů je Veličenstvo Kat. :/
   	
 \refren  	
	Nad palácem vlády ční prapor s gilotinou,
 
 	děti mají rády kornouty se zmrzlinou,
  
  	soudcové se na ně zlobili,
   
   	zmrzlináře dětem zabili.


\end{minipage}\begin{minipage}[t]{0.55\textwidth}\setlength{\parindent}{0.45cm}\vspace*{0.55cm}  % V případě varianty č.2 jde odsud text do pravé části

\sloka
	Byl hrozný tento stát, když musel jsi se dívat,
   
   	jak zakázali psát a zakázali zpívat,
   
   	a bylo jim to málo, poručili dětem
   
   	/: modlit se, jak si přálo Veličenstvo Kat. :/

\sloka
	S úšklebkem Ďábel viděl pro každého podíl,
  
  	syn otce nenáviděl, bratr bratru škodil,
   
   	jen motýl smrtihlav se nad tou zemí vznáší,
   
   	/: kde v kruhu tupých hlav dlí Veličenstvo Kat. :/


\end{minipage}
}
\setcounter{Slokočet}{0}
\end{song}


\end{document}
